\chapter*{Annexe}\addcontentsline{toc}{chapter}{Annexe}\label{annexe}

\section*{Exemple de requête API}
Le fichier \texttt{body.json} situé à la racine du projet fournit un exemple de requête pour l'API de diagnostic.

\begin{lstlisting}[language=JSON, caption={Exemple de payload JSON}, label={lst:bodyjson}]
{
  "title": "Coupure Internet fibre",
  "description": "Client acc-12345, service svc-fiber-12345, commande ord-2026-001. Coupure Internet fibre.",
  "source": "chat",
  "priority": "P3"
}
\end{lstlisting}

\section*{Variables d'environnement}
Le tableau ci-dessous reprend les principales variables d'environnement du fichier \texttt{.env.example}.

\begin{table}[H]
    \centering
    \renewcommand{\arraystretch}{1.2}
    \begin{tabular}{|l|l|}
        \hline
        \rowcolor{iNavy} \textbf{\color{white}Variable} & \textbf{\color{white}Valeur par défaut} \\ \hline
        MOCK\_LLM & true \\ \hline
        OLLAMA\_BASE\_URL & http://localhost:11434 \\ \hline
        OLLAMA\_MODEL & qwen2.5:latest \\ \hline
        NEO4J\_URI & bolt://localhost:7687 \\ \hline
        NEO4J\_USER & neo4j \\ \hline
        NEO4J\_PASSWORD & password \\ \hline
        EMBEDDING\_PROVIDER & ollama \\ \hline
        EMBEDDING\_MODEL & mxbai-embed-large:latest \\ \hline
        VECTOR\_INDEX\_DIM & 1024 \\ \hline
        VECTOR\_SIMILARITY\_THRESHOLD & 0.75 \\ \hline
    \end{tabular}
    \caption{Extrait des variables d'environnement}
\end{table}

\section*{Commandes utiles}
\begin{lstlisting}[language=bash, caption={Commandes principales du Makefile}]
make install      # pip install -r requirements.txt
make seed         # Seed Neo4j local (schema + mocks + embeddings)
make test         # pytest tests/  (Neo4j local requis)
make api          # uvicorn api.main:app --reload
make docker-up    # docker compose up --build
make docker-down  # docker compose down
\end{lstlisting}

\section*{Requête cURL}
Pour tester l'API en mode texte simplifié :

\begin{lstlisting}[language=bash, caption={Appel de l'API diagnose/text}]
curl -X POST "http://127.0.0.1:8000/diagnose/text" \\
  -H "Content-Type: application/json" \\
  -d @body.json
\end{lstlisting}

\section*{Structure du projet}
\begin{lstlisting}[language=bash, caption={Arborescence simplifiee}]
diagnostic-technique/
├── api/
├── config/
├── graph/
│   ├── schema.cypher
│   ├── queries.py
│   └── ingestion/
├── orchestrator/
│   └── pipeline.py
├── sub_agents/
│   ├── intake_parser/
│   ├── logs_investigator/
│   ├── context_agent/
│   ├── rag_ticket_search/
│   ├── root_cause_reasoner/
│   ├── ticket_manager/
│   ├── remediation_explainer/
│   ├── content_guardrail/
│   ├── report_generator/
│   └── notifier/
├── tests/
├── ui/
├── docker/
└── Makefile
\end{lstlisting}
